% Generated draft: requires mathematical and source review.
\documentclass[11pt]{article}
\usepackage[margin=25mm]{geometry}
\usepackage{fontspec}
\setmainfont{Latin Modern Roman}
\usepackage{amsmath,amssymb,mathtools,unicode-math}
\setmathfont{Latin Modern Math}
\usepackage{enumitem,needspace,xurl,hyperref}
\setlength{\parindent}{0pt}
\setlength{\parskip}{5pt}
\setlength{\emergencystretch}{3em}
\newcommand{\N}{\mathbb{N}}
\newcommand{\Z}{\mathbb{Z}}
\newcommand{\Q}{\mathbb{Q}}
\newcommand{\R}{\mathbb{R}}
\newcommand{\C}{\mathbb{C}}
\newcommand{\F}{\mathbb{F}}
\newcommand{\A}{\mathbb{A}}
\newcommand{\PP}{\mathbb{P}}
\newcommand{\E}{\mathbb{E}}
\newcommand{\eps}{\varepsilon}
\newcommand{\dd}{\,\mathrm{d}}
\begin{document}
\section*{UnsolvedMath ID: 701041}
\textbf{Source problem code:} catalogue entry
\subsection{Definitions and mathematical statement}
Let \(p\) be a prime. A pro-\(p\) group is a topological group that is an inverse limit of finite \(p\)-groups. For a pro-\(p\) group \(H\), define its closed lower central series by
\[
\gamma_1(H)=H,\qquad
\gamma_{n+1}(H)=\overline{\langle [x,y]:x\in\gamma_n(H),\ y\in H\rangle}\quad(n\geq 1),
\]
where \([x,y]=x^{-1}y^{-1}xy\) and the closure is taken in \(H\).

For a finite field \(k\) of characteristic \(p\), the Nottingham group over \(k\) is
\[
J(k)=\left\{t+\sum_{n\geq 2}a_nt^n:a_n\in k\right\},
\]
with formal composition as its group operation and the topology whose identity neighborhoods are
\[
U_m=\{f\in J(k):f(t)\equiv t\pmod{t^{m+1}}\}\qquad(m\geq 1).
\]
The supplied record denotes the Nottingham group by \(J\) but does not specify its coefficient field. Fixing the intended field \(k\) and writing \(J=J(k)\), the requested assertion is:
\begin{quote}
For every infinite pro-\(p\) group \(G\), if
\[
G/\gamma_{2p+1}(G)\cong J/\gamma_{2p+1}(J),
\]
then \(G\cong J\).
\end{quote}
Here quotient groups carry their quotient topologies, and \(\cong\) is interpreted as an isomorphism of topological groups. The use of closed lower central series and topological isomorphisms is the pro-\(p\) interpretation adopted for this draft; the record does not explicitly specify these conventions. The coefficient field remains unresolved: the displayed formulation is conditional on its intended choice, rather than an assertion supplied for every finite field of characteristic \(p\). The record states no restriction on the prime \(p\).
\subsection{Short English statement}
The proposed assertion says that an infinite pro-\(p\) group is determined, up to topological isomorphism with the intended Nottingham group \(J\), by its quotient modulo \(\gamma_{2p+1}\). The coefficient field defining \(J\) is unspecified in the supplied record and requires clarification.
\subsection{Sources}
\begin{itemize}
\item Supplied local UnsolvedMath registry record: permanent ID \(701041\), exact source code \texttt{catalogue entry}, title \emph{Kourovka Notebook Problem 14.56}; upstream revision \texttt{d371ee917ac9647fbb06cd6fc529deec6f3dda82}. The supplied status is \texttt{open}; this label is not an independently verified mathematical or literature assessment.
\item Attributed in the supplied record to C. R. Leedham-Green, proposed in 1999, in \emph{The Kourovka Notebook (Unsolved Problems in Group Theory)}, 14th issue, Problem 14.56.
\item The supplied background identifies the 21st issue edition (2026), arXiv:1401.0300v48: \url{https://arxiv.org/abs/1401.0300v48}. The supplied collection description names E. I. Khukhro and V. D. Mazurov as editors. These bibliographic details are supplied attribution; the cited source has not been independently inspected for this draft.
\end{itemize}
\end{document}
